\section{Correctness Arguments}
\label{sec-proof}

We briefly establish key properties of the selective view materialization algorithm, including the minimal form uniqueness, the termination guarantee, and the completeness, which together form the correctness of \proto algorithm.
Due to limited space, we leave the full proof and details in the supplementary appendix ~\cite{proof}.

\medskip \noindent {\bf Minimal Form Uniqueness.}
We first define the minimal form of a materialization plan, and then prove the uniqueness of the minimal form for plans generated by our algorithm.
%We denote a valid view materialization plan as a valid set for conciseness.

\begin{deff}
A relation set $S'$ is a minimal form of another relation set $S$ if and only if:
\begin{enumerate}
    \item 
    $S' \subseteq S$, 
    \item
    $S'$ is a minimal view materialization plan (Definition~\ref{def:min-plan}).
\end{enumerate}
\end{deff}

\begin{lmm}\label{lemma:uniqness}
    A valid materialization plan $S$ generated following the base minimal plan has a unique minimal form $S'$.
\end{lmm}

\begin{proof}
    We proceed by establishing contradictions. 
    Assume a valid materialization plan has two distinct minimal forms X and Y.

    If $X \subset Y$ (or vice versa), we can remove extra relations in Y to get a smaller valid set, leading to a contradiction.
    
    Otherwise, let $A=X \setminus Y$, $B= Y \setminus X$.
    We claim that each relation $a \in A$ depends on at least one relation in B, while each relation $b \in B$ depends on some relation in A, which creates a contradictory circular dependency.
    To see this, relational queries to $a$ should be executable given $Y$.
    However, if $a$ is executable solely on $X \cap Y$, we can obtain another valid set by removing $a$ from $X$,
    which contradicts that $X$ is in a minimal form.
\end{proof}

\medskip\noindent{\bf Termination.} The enumeration algorithm (Algorithm~\ref{alg:enumerate}) traverses the replacement graph. Given that both the contract and the replacement graph should have no cycle, the traversal is guaranteed to terminate at nodes without any outgoing edges.
    

\medskip \noindent{\bf Completeness.}
We first define the $Replace$ and $GetMinimal$ functions and then establish the completeness of our algorithm.

Given a set of view materialization plans $\Sigma$, we define:
\[ Replace(\Sigma) \triangleq \{ Replace(S) ~|~ S \in \Sigma \} \]
\[ GetMinimal(\Sigma) \triangleq \{ GetMinimal(S) ~|~ S \in \Sigma \} \]

Let $B$ denote the base materialization plan in Algorithm~\ref{alg:enumerate} and $B'$ denotes its minimal form.
\begin{deff}
We define {\em all possible view materialization plans} obtained by applying the Replace function for $k$ times as:
       \[ \Sigma_k \triangleq \bigcup_{i=0}^{k} Replace^i(B) \] 
\end{deff}

\begin{deff}\label{def:min-replace}
We define $MinReplace$ as the combination of $Replace$ and $GetMinimal$:
        \[ MinReplace(\Sigma) \triangleq GetMinimal(Replace(\Sigma))  \]
\end{deff}
        
With this definition, the output of Algorithm~\ref{alg:enumerate} is effectively:
\[  \bigcup_{i=0}^{K}MinReplace^i(B') \]

\begin{lmm}\label{lmm:getminimal}
For any valid materialization plan $S$, we claim:
\[
    MinReplace(\{S\}) = MinReplace(GetMinimal(\{S\}))
\]
\end{lmm}

This lemma can be further developed to fit a set of valid materialization plans,
instead of just one plan.

\begin{thm}\label{theorem:completeness}{Completeness.}
Algorithm~\ref{alg:enumerate} outputs all minimal forms of view materialization plans obtainable from $\Sigma_K$:
        \[ \bigcup_{i=0}^{K}MinReplace^i(B') =  GetMinimal(\Sigma_K) \]
\end{thm}

This theorem asserts that the application of the $GetMinimal$ function during each iteration of the algorithm does not compromise the final correctness of the algorithm. Instead, it enhances efficiency by early pruning of non-minimal materialization plans.
    

Given Lemma~\ref {lmm:getminimal}, we can prove completeness by induction on the $MinReplace$ step. 
Essentially, at each iteration $k$, the application of $MinReplace$ function maintains completeness regarding the view materialization plan, 
considering $k$ applications of the $replace$ function on the base materialization plan $\Sigma_k$.
The complete inductive steps are provided in the supplementary material ~\cite{proof}.

